\documentclass[10pt,a4paper]{amsart}
\usepackage[margin=19mm]{geometry}
\usepackage{amsmath,amssymb,mathtools}
\usepackage[scr=rsfs,cal=euler]{mathalfa}
\usepackage{xfrac}
\usepackage[unicode,hidelinks,bookmarks=false]{hyperref}
\newcommand{\ie}{i.e.\ }
\newcommand{\cf}{cf.\ }
\newcommand{\resp}{respectively\ }
\newcommand{\Subsection}{\subsection}
\providecommand{\nobreakdash}{\nobreak\hbox{-}}
\setlength{\emergencystretch}{2em}
\multlinegap0pt
\usepackage{bm}
\usepackage{bbm}
\usepackage{dsfont}
\usepackage{xspace}
\usepackage{cancel}
\usepackage{mathtools}
\usepackage{amsthm}
\makeatletter
\@ifundefined{Theorem}{\newtheorem{Theorem}{Theorem}}{}
\@ifundefined{Lemma}{\newtheorem{Lemma}[Theorem]{Lemma}}{}
\@ifundefined{Remark}{\newtheorem{Remark}[Theorem]{Remark}}{}
\makeatother
\makeatletter
\def\R{\mathbb R}
\newcommand{\eps}{\varepsilon}
\newcommand{\gm}{g_{\mathrm{Min}}}
\newcommand{\mR}{\mathbb{R}}                    % Formatting for R
\def\mid{\,:\,}
\expandafter\ifx\csname tehtratk\endcsname\relax
\newenvironment{tehtratk}%
             {\begin{list}{\arabic{enumi}.}{\usecounter{enumi}%
              \setlength{\labelsep}{0.5em}%
              \settowidth{\labelwidth}{(\arabic{enumi})}%
              \setlength{\leftmargin}{\labelwidth+\labelsep}}}%
\fi
\makeatother
\title[Semiglobal uniqueness for the Lorentzian Calder\'on problem]{Semiglobal uniqueness for the Lorentzian Calder\'on problem}
\author{Lauri Oksanen, Rakesh, Mikko Salo}
\date{Source: arXiv:2608.13116v1 (CC BY 4.0)}
\begin{document}
\maketitle

\noindent\emph{Source and licence.} Semiglobal uniqueness for the Lorentzian Calder\'on problem. Version arXiv:2608.13116v1 (CC BY 4.0), distributed under the Creative Commons Attribution 4.0 International licence, as recorded in the arXiv licence metadata for this exact version and shown on the abstract page https://arxiv.org/abs/2608.13116v1. Full rights notice: licenses/arxiv-2608.13116-CC-BY-4.0.txt.\par\smallskip

\noindent\textit{Editorial scope.} Counted unit: the Lemma whose source label is \texttt{lem\_nontrapping} in the article (source0.tex lines 1097--1104, source label \texttt{lem\_nontrapping}), delivered together with the article's own complete original proof of it (source0.tex lines 1105--1161, one \texttt{proof} environment of 57 lines). No mathematics is written, repaired or re-derived: statement and proof are the article's own text line for line. Required context retained uncounted: assumption1 (lines 231--234); assumption2 (lines 231--234); lem\_r\_plus (lines 999--1006); rem\_alpha\_diffeo (lines 1053--1055); eq\_beta (lines 1088--1092). Every definition, display and label of those lines is retained unchanged. Omitted from this package and not redistributed: 1--230, 235--998, 1007--1052, 1056--1087, 1093--1096, 1162--2721. Nothing inside a retained excerpt is omitted or altered: every retained body is delivered exactly as the article's lines are written, so no omission receipt of any kind accompanies this package. Citations: the retained excerpts carry no citation command at all, so no bibliography entry of the article is transcribed and no reference is redistributed. Reference adaptation: none was needed and none was made; every reference inside the delivered unit resolves inside it, so no unresolved reference or citation is produced. Printed slips kept literal: no printed word, symbol, hypothesis or number of the counted statement or of its proof was altered. Layout and class: the article's own class is \texttt{amsart} with packages including amsmath, amscd, amssymb, amsthm, mathtools, graphicx, mathrsfs, hyperref, calc, comment, tikz, xcolor; the wrapper is the lane's pinned \texttt{amsart} editorial wrapper at 10pt on A4, which restates the theorem-like environments the retained text needs and adds the article's own macro definitions that the retained text needs (\texttt{\textbackslash R}, \texttt{\textbackslash eps}, \texttt{\textbackslash gm}, \texttt{\textbackslash mR}, \texttt{\textbackslash mid}), with extra wrapper packages (bm, bbm, dsfont, xspace, cancel, mathtools). The delivered numbering is the unit's own. Assets: no retained span contains a figure environment, a graphics-inclusion command, a PDF-page inclusion, a table or a TikZ picture; no image file is copied into this package, the package declares no asset dependency and no image file is redistributed with it. Licence: version arXiv:2608.13116v1 of this article is distributed under the Creative Commons Attribution 4.0 International licence, as recorded in the arXiv licence metadata for this exact version and shown on the abstract page https://arxiv.org/abs/2608.13116v1. A full rights notice is delivered at licenses/arxiv-2608.13116-CC-BY-4.0.txt. Characters: every retained excerpt is pure ASCII, so the delivered engine, XeLaTeX with its default Latin Modern text font, reports no missing character.
\begin{gather}
g = \gm \text{ outside $M$}, \label{assumption1} \\
g \text{ is globally hyperbolic in $\mR^{1+n}$.} \label{assumption2}
\end{gather}

\begin{Lemma}\label{lem_r_plus}
Suppose that $g$ satisfying \eqref{assumption1}--\eqref{assumption2} is simple in direction $\omega \in S^{n-1}$. Then there is a smooth function $r_+ : \Sigma_- \to (0, \infty)$ and a constant $\eps > 0$ such that $r_+(z) = 2 + \eps$ when $|z|$ is large and, writing
    \begin{align*}
\mho_\delta &= \{(z, r) \mid r \in (-\infty, r_+(z) + \delta),\ z \in \Sigma_- \}, \quad \delta \in \R, 
    \end{align*}
it holds that $\R \times B \subset \Phi^g(\mho_{-\eps})$ and that
$\Phi^g : \mho_\eps \to \Phi^g(\mho_\eps)$ is a diffeomorphism.
\end{Lemma}

\begin{Remark}\label{rem_alpha_diffeo}
If $g$ satisfying \eqref{assumption1}--\eqref{assumption2} is simple in direction $\omega$, then in the coordinates given by $\Phi^g$, the geodesic $\gamma_z$, $z \in \Sigma_-$, is the curve $r \mapsto (z, r)$. In particular, $\gamma_z$ is transverse to $\Sigma_+^g$, and the map $\alpha(z) = \gamma_z(r_+(z))$ is a diffeomorphism from $\Sigma_-$ to $\Sigma_+^g$.
\end{Remark}

    \begin{align}\label{eq_beta}
\beta_{g}(\Sigma_{-,\sigma} \setminus \mathcal T_{g})
= 
\beta_{g'}^{\lambda}(\Sigma_{-,\sigma}),
    \end{align}

\begin{Lemma}\label{lem_nontrapping}
Let two Lorentzian metrics $g$ and $g'$ on $\R^{1+n}$ both satisfy \eqref{assumption1} and \eqref{assumption2}. 
Suppose that $g'$ is simple in direction $\omega \in S^{n-1}$ and that there is $\lambda : \Sigma_+^{g'} \to \R$ such that \eqref{eq_beta} holds for all $\sigma \in \R$.
Then 
$\mathcal{T}_{g} = \emptyset$, $\gamma_z^{g}$ intersects $\Sigma_+^{g'}$ transversally for all $z \in \Sigma_-$,
$r_+^{g}$ is smooth on $\Sigma_-$, the map
taking $z$ to $\gamma_z^{g}(r_+^{g}(z))$ is a diffeomorphism from $\Sigma_{-}$ to $\Sigma_{+}^{g'}$, and $\lambda$ is smooth and positive.
\end{Lemma}

\begin{proof}
Since $g$ and $g'$ are globally hyperbolic, they are time-orientable, and we fix their time-orientations so that $dt$ is future-directed outside $M$ for both the metrics.
Let $w \in \Sigma_+^{g'}$.
Using the directional simplicity of $g'$, Remark \ref{rem_alpha_diffeo} implies that
    \begin{align}\label{def_line_bundle}
\{(\gamma_z^{g'}(r), \lambda_z \dot \gamma_z^{g'}(r)) \mid r=r_+^{g'}(z),\ \lambda_z \ne 0, \ z \in \Sigma_{-}\}
    \end{align}
is a line bundle over $\Sigma_{+}^{g'}$.
In particular, there is unique $z' \in \Sigma_-$ such that $w = \gamma_{z'}^{g'}(r_+^{g'}(z'))$. It follows from \eqref{eq_beta} that there
are $z \in \Sigma_{-}$, $s > 0$ and $\lambda \in \R$ such that $\gamma_z^{g}(s) = w$ and 
    \begin{align}\label{def_lambda}
\dot \gamma_z^{g}(s) = \lambda \xi, \quad \xi = \dot \gamma_z^{g'}(r_+^{g'}(z')).
    \end{align}
The geodesics $\gamma_z^{g}$ and $\gamma_{z'}^{g'}$ are future pointing at $z$ and $z'$, respectively. Thus they are future pointing everywhere, and therefore $\lambda > 0$. 

Consider the largest number $\ell \in (0, s]$ such that $\gamma_z^{g}(r)$ is outside the cylinder $\R \times B$ for $r \in (s-\ell, s)$. Then $\gamma_z^{g}|_{(s-\ell, s)}$ coincides with a segment of $\gamma_{z'}^{g'}$, up to reparametrization. This segment of $\gamma_{z'}^{g'}$ does not intersect $\Sigma_+^{g'}$ because of Lemma \ref{lem_r_plus}. Moreover, if $\gamma_z^{g}$ enters into the cylinder $\R \times B$, then it can not intersect $\Sigma_+^{g'}$ before it exits. Since $B$ is convex, it follows from \eqref{assumption1} that if $\gamma_z^g$ exits the cylinder then it does not re-enter. We conclude that $s = r_+^{g}(z)$.

The geodesics through $(w, \lambda \xi)$, for $\lambda \ne 0$, coincide up to reparametrization. This implies that $z \in \Sigma_{-}$ is unique, and we write $\kappa(w) = z$. Then $w = \alpha(\kappa(w))$ where $\alpha(z) = \gamma_z^{g}(r_+^{g}(z))$. By \eqref{def_lambda} and Remark \ref{rem_alpha_diffeo} the geodesic $\gamma_z^{g}$ is transverse to $\Sigma_+^{g'}$ for $z \in \kappa(\Sigma_+^{g'})$, and therefore $r_+^{g}$ is smooth near any $z \in \kappa(\Sigma_+^{g'})$. Hence $\alpha$ is smooth in $\kappa(\Sigma_+^{g'}) \subset \Sigma_-$. 

Write $|\cdot|$ for the Euclidean norm and $\alpha'(z) = \gamma_z^{g'}(r_+^{g'}(z))$. Let us show that $\kappa$ is smooth. 
Using $\lambda > 0$, the function 
    \begin{align*}
\xi(w) 
= 
\frac{\dot \gamma_{\kappa(w)}^{g}(r_+^{g}(\kappa(w)))}
{|\dot \gamma_{\kappa(w)}^{g}(r_+^{g}(\kappa(w)))|} 
= 
\frac{\dot \gamma^{g'}_{z'}(r_+(z'))}
{|\dot \gamma^{g'}_{z'}(r_+(z'))|},
\quad 
z' = (\alpha')^{-1}(w),
    \end{align*}
is smooth. Define
    \begin{align*}
r_-(z) = \inf \,\{ r > 0 \,:\, \gamma^{g}(-r; w, \xi(w)) \in \Sigma_- \}.
    \end{align*}
Then
    \begin{align*}
\kappa(w) = \gamma^{g}(-r_-(w); w,\xi(w)).
    \end{align*}
Moreover, $\gamma^{g}(\cdot; w,\xi(w))$ is transverse to $\Sigma_-$ since
    \begin{align*}
\dot \gamma(-r_-(w); w,\xi(w))
=
\tilde \lambda \dot \gamma_{\kappa(w)}^{g}(0)
    \end{align*}
for some $\tilde \lambda > 0$.
The transversality implies that $r_-$ is smooth. Therefore $\kappa$ is smooth.

Since $\alpha$ is a smooth left inverse of $\kappa$, $\kappa$ is a local diffeomorphism. In the coordinates given by $\Phi^{g'}$, we can identify $z$ and $(z, 0)$. Then it holds for $z \in \Sigma_-$ with large $|z|$ that $\kappa(z, 2) = z$. Hence $\kappa$ is proper. It follows from Hadamard's global inverse function theorem that $\kappa$ is a diffeomorphism from $\Sigma_+^{g'}$ to $\Sigma_-$. This again implies that $\mathcal T_{g} = \emptyset$ and that $\alpha$ is a diffeomorphism from $\Sigma_-$ to $\Sigma_+^{g'}$.

Finally, the smoothness of $\lambda$ follows from 
    \begin{align*}
\lambda(w) = |\lambda(w)| = \frac{|\dot \gamma_{z'}^{g'}(r_+^{g'}(z'))|}{|\dot \gamma_z^g(r_+^{g}(z))|},
    \end{align*}
where $z = \alpha^{-1}(w)$ and $z' = (\alpha')^{-1}(w)$.
\end{proof}

\end{document}
